% Draft, 31-10-2026. Every number below is a measurement in ../README.md or in the training logs it cites (same laptop, same
% code); the related work was checked against the papers themselves on 31-10-2026.
% Build: ./build.sh (docker texlive, here or where oz_tex of ../.hosts runs)
\documentclass[11pt]{article}
\usepackage[T1]{fontenc}
\usepackage{lmodern}
\usepackage[a4paper,margin=2.4cm]{geometry}
\usepackage{amsmath,amssymb}
\usepackage{booktabs}
\usepackage{hyperref}
\usepackage{xcolor}
\hypersetup{colorlinks=true,linkcolor=blue!50!black,citecolor=blue!50!black,urlcolor=blue!50!black}

\title{Every product correctly rounded:\\bit-reproducible LoRA training of a 375-billion-parameter\\mixture of experts from a laptop's disk}
\author{Shaban Ejupi\\Independent researcher}
\date{31 October 2026}

\begin{document}
\maketitle

\begin{abstract}
Floating-point training is not a function of its inputs. GPU kernels choose their reduction order by shape, backward passes
accumulate with atomics, and the same data in another batch, on another GPU or with another library gives other gradients.
Recent batch-invariant and deterministic kernels fix one order per kernel: a run becomes repeatable on the same software and
hardware, but its result is still defined by an implementation. We take the other route: every linear map of the forward and
of the backward pass returns the \emph{correctly rounded} value of its exact sum, a number defined by the mathematics alone.
The products run exactly on INT8 tensor cores from digit planes of BF16 and FP8 operands; the rounding is certified, and every
output the certificate leaves open is recomputed exactly on the CUDA cores or, as a last resort, in a 2560-bit host
superaccumulator. Gradients through FP8 expert weights are exact, never quantized. Applied to rank-16 LoRA on
Jais-3-375B-A23B-FP8 (61 layers, 192 FP8 experts, a 391~GB checkpoint) streamed from the NVMe disk of a laptop with a 12~GB GPU
once per forward and once per backward, a step of $8\times1024$ tokens takes about 30 minutes (1\,792--2\,009~s). The step does
not depend on how rows are grouped into products (4\,096 or 40 rows), on the number of threads (24 or 7), on interruption (a
resumed run is byte-identical) or on which of three generations of our product kernels computes it; no output is ever used
uncertified. Five hand-written backward formulas agree with central differences of their FP64 forward to
$5.1\times10^{-7}$, and finite differences through the model and the predicted loss change of each update agree with the
backward. Trained on 191.9 million tokens of public documents of Kosovo's institutions, the held-out loss falls from 1.5102
at step~1 to 1.463--1.468 at steps 15--18.
\end{abstract}

\noindent\textbf{Përmbledhje (shqip).} Trajnimi me numra me presje lëvizëse nuk është funksion i hyrjeve të tij: e njëjta
e dhënë në një grup tjetër, në një GPU tjetër ose me një librari tjetër jep gradientë të tjerë. Këtu çdo shumëzim matricash i
trajnimit, përpara dhe prapa, jep vlerën e \emph{rrumbullakosur saktë} të shumës së saktë---një numër i përcaktuar vetëm nga
matematika---i llogaritur mbi tensor core-t INT8, i certifikuar, dhe i rillogaritur saktë kur certifikata s'mjafton. Me këtë
trajnuam LoRA me rang 16 mbi Jais-3-375B-A23B-FP8 (61 shtresa, 192 ekspertë FP8, 391~GB) nga disku NVMe i një laptopi me
GPU 12~GB: një hap me $8\times1024$ tokena zgjat rreth 30 minuta, dhe nuk varet nga grupimi i rreshtave, numri i fijeve,
ndërprerja apo versioni i kernel-ave; asnjë dalje e pacertifikuar nuk përdoret kurrë. Mbi 191,9 milionë tokena dokumentesh
publike të institucioneve të Kosovës, humbja e mbajtur jashtë bie nga 1,5102 në 1,463--1,468.

\section{Introduction}

Two runs of the same training script on the same data rarely produce the same model. Floating-point addition is not
associative, and a GPU kernel's order of summation depends on the shape of the problem (how $K$ is split across thread
blocks), on the library version and on the hardware; backward passes add partial gradients with atomic instructions whose
order is decided at run time. Zhuang et al.\ measured the consequence: nondeterministic tooling barely moves aggregate accuracy
but makes performance on subsets of the data much more variable, and enforcing determinism cost up to 746\%
overhead~\cite{zhuang21}.

The current answer is to fix the order. Batch-invariant kernels make a token's output independent of the batch it is computed
in~\cite{he25}; tree-based invariant kernels align the reduction order across tensor-parallel sizes~\cite{tbik};
fixed-configuration GEMMs give the same bits across three NVIDIA architectures for inference~\cite{reprofuse}; on-policy
reinforcement learning with batch-invariant forward kernels and hand-written backward passes keeps trainer and sampler bitwise
equal~\cite{vllm25}; and DeepSeek-V4 replaces cuBLAS end to end, abandons split-$K$ and replaces atomics in the attention and
MoE backward by per-SM buffers and deterministic reductions, for training that is repeatable and aligned with
inference~\cite{dsv4}. All of these define the result by an implementation: the bits are those of one kernel with one tiling.
Another kernel, another tile shape or another GPU generation gives other bits, and nobody outside can check a gradient without
running the same code on the same kind of machine.

We define the result by the mathematics instead. Every linear map of the training step---every projection, every expert, the
output layer, in the forward \emph{and} the backward pass---returns the correctly rounded value of its exact sum. That value is
unique, so it cannot depend on the order of summation, the tiling, the batch, the thread count or the machine; any third party
can recompute any output with exact rational arithmetic and compare. INT8 tensor cores make this affordable: an INT8 product
with INT32 accumulation is exact, and a BF16 or FP8 operand is a short sum of INT8 digits times powers of two, so the exact
product is a sum of exact INT8 products~\cite{ozaki2,companion}. What remains is to know when the rounding is decided, and to
finish the rare outputs where it is not.

The test case is deliberately out of proportion to the hardware. Jais-3-375B-A23B-FP8 (379 billion parameters, Apache-2.0,
61 layers, 192 routed FP8 experts per MoE layer) does not fit in the memory of any single GPU; we fine-tune it with rank-16 LoRA
adapters~\cite{hu21} on a laptop whose GPU (RTX PRO 3000 Blackwell Laptop, sm\_120, 12~GB) holds one layer at a time,
streaming the checkpoint from the laptop's NVMe disk twice per step. Our contributions:
\begin{itemize}
\item correctly rounded products for every operand format of a modern FP8 MoE in training: BF16 projections with BF16 or FP32
outputs, FP8 experts with block scales as the model prescribes (W8A8), and the transposed products of the backward, including
an exact gradient through FP8 weights (Section~\ref{sec:products});
\item a complete LoRA training step of a 375B MoE from disk on one laptop GPU, with every host-side operation in a fixed FP64
order (Section~\ref{sec:step});
\item evidence that the step is a function of its inputs: byte-identical gradients across rows per product, thread counts,
interruption and three generations of product kernels, and gradient checks (Section~\ref{sec:verify});
\item measurements of what exactness costs and where the time goes, and a first training run on public data
(Sections~\ref{sec:perf}, \ref{sec:train}).
\end{itemize}
To our knowledge, no earlier training run has had every product of its forward and backward pass correctly rounded, and no
earlier LoRA fine-tuning of a model of this size has run on a single laptop GPU.

\section{Background}
\paragraph{LoRA.} A frozen map $y = W x$ gains a trainable low-rank term, $y = W x + s\, B A x$ with $A\in\mathbb{R}^{r\times
k}$, $B\in\mathbb{R}^{m\times r}$ and scale $s=\alpha/r$~\cite{hu21}. We follow the PEFT library's defaults: $B=0$ and $A$
uniform in $\pm1/\sqrt{k}$, so the first step's $\partial L/\partial A$ is exactly zero. QLoRA~\cite{qlora} quantizes the
frozen weights to fit a GPU; offloading systems move parameters and optimizer state to CPU memory and NVMe instead, and
fine-tune trillion-parameter models on one DGX-2 node (16 GPUs)~\cite{zeroinf} or reduce the host memory that SSD offloading
needs~\cite{memascend}. We keep the checkpoint's own formats and stream it, one layer at a time, through a single GPU.

\paragraph{Fine-grained FP8.} Jais-3's experts store E4M3 weights with one BF16 scale per $128\times128$ block and quantize
each token's activation per group of 128 with an FP32 scale before the product (the W8A8 scheme DeepSeek-V3 introduced at
scale~\cite{dsv3}). Attention projections, the shared expert, the router and the output layer are BF16.

\paragraph{Exact products on integer units.} The Ozaki scheme splits floating-point matrices into integer slices whose
products integer tensor cores compute exactly~\cite{ozaki2}; our companion paper builds FP64, multi-word and arbitrary
precision on the INT8 units of the same laptop GPU and correctly rounded MXFP4, NVFP4 and BF16 products for inference of the
largest open models~\cite{companion}. Reproducible and correctly rounded summation on CPUs and GPUs has a long history---long
accumulators~\cite{kulisch}, reproducible BLAS~\cite{reproblas,exblas} and accurately rounded BLAS~\cite{rareblas}---but
not, as far as we know, inside a training loop.

\section{Correctly rounded products for training}\label{sec:products}

\paragraph{BF16 $\times$ BF16.} For $Y=XW^{T}$ each block of 128 consecutive $k$ of a weight row or a token row shares a grid at
its largest exponent; on that grid each element is an integer, cut into signed 7-bit digits. Weights take 3 digits and tokens 4
when the output is BF16 ($p=8$), 5 and 6 when it is FP32 ($p=24$). Each pair of digit planes is one INT8 GEMM with INT32
accumulation (\texttt{mma.m16n8k32}), and the partial sums of a block are combined in an \texttt{int64} window at the token's
largest block, with headroom for the number of blocks. Elements below the grid are floored; their count bounds the tail, and
the output is certified when the whole interval rounds to one floating-point value (round to nearest even, subnormals
included). Our tiled kernel computes 64 rows $\times$ 32 tokens per block of 256 threads, streaming $K$ through 3--4
\texttt{cp.async} stages and loading fragments with \texttt{ldmatrix}.

\paragraph{The fixup.} An uncertified output is written as NaN and recomputed exactly on the CUDA cores, one warp per output:
every product is formed exactly and summed in a 128-bit fixed point anchored at the largest product of that row and token. If
even that window cannot certify it, the output is finished on the host in a 2\,560-bit superaccumulator with no window and no
tail, which always succeeds. Nothing uncertified is ever used.

\paragraph{FP8 experts, forward.} The activation is first quantized exactly as the model prescribes (E4M3 per group of 128 with
an FP32 scale); the product of E4M3 weights and E4M3 activations is then exact in three digits each, and the block term
$(S\,s)\sum w x$ fits a $64\times32\to128$-bit product shifted into the 63-bit window. Outputs the window leaves open are
recomputed exactly as the reference FP8 kernel does, so every output---NaN or not---equals that kernel's.

\paragraph{The backward through FP8 weights, exactly.} The gradient of an input is $dX = dY\,W$. We round $dY$ once to BF16,
the gradient format of BF16 training, and never to FP8. An E4M3 weight (4 significant bits) times its BF16 block scale
(8 bits) is exact in FP32 and splits exactly into two BF16 numbers $w_{hi}+w_{lo}$, so $[\,dY\,|\,dY\,]\,[W_{hi};W_{lo}]$ is the
exact product with the FP8 weights. Our current kernel avoids the doubling: it transposes the weight's own E4M3 digits (three,
exact) and applies the block scale per block of 128, so a term costs $3\times6=18$ digit products instead of
$2\times4\times6=48$, with the same correctly rounded outputs. Maps that share an input form one product (q, k, v together;
an expert's gate and up together), rounded once.

\paragraph{LoRA gradients.} With $G=\partial L/\partial(\text{projection output})$ before RoPE and $H=\text{BF16}(x A^{T})$,
$\partial L/\partial B = s\,G^{T} H$ is one certified product over the tokens ($K$ = number of tokens). No bit of it can depend
on the order of the tokens or on how rows were cut into products.

\section{A training step of a 375B MoE on a laptop}\label{sec:step}

\paragraph{Forward.} The model runs layer by layer over all rows of the step (8 training sequences of 1\,024 tokens and 2
held-out sequences), each layer read from disk once and applied to every row; the residual stream entering each layer is kept
in host memory (activation checkpointing at layer granularity~\cite{chen16}). Each expert is read once per layer and applied to
all rows routed to it, in ascending expert id, and each row's weighted sum of its experts is accumulated in FP64 in that order
and rounded once. Run this way, the batched model gives the token-by-token model bit for bit~\cite{companion}.

\paragraph{Backward.} From the top, every layer is read from disk once more and recomputed from its saved input with the
forward's own functions; then every map's $dX$ is one correctly rounded product (FP32 outputs) added to an FP64 gradient.
Attention runs in FP64 with every sum in a fixed order: the scores, the probabilities' gradient, $dq$, $dk$ and $dv$ on the
GPU's FP64 units with \texttt{\_\_dmul\_rn}/\texttt{\_\_dadd\_rn} in exactly the host's order, the softmax with the host's
\texttt{exp}. RoPE's inverse rotation, RMSNorm and SwiGLU are FP64 on the host. The routing weights $w_e = 2.5\,s_e/\sum s$
(sigmoids of the eight chosen experts) pass their gradient through the sigmoid to the router logits and its transposed product;
the top-8 choice is a constant, as in every MoE backward. The backward reads only the experts the training rows were routed
to: 376--379~GB per step.

\paragraph{Streaming.} A reader thread fetches the tensors in the order they will be used into a pinned ring; the time spent
waiting for it counts as disk time, and in the backward it is 0.0~s per step: the disk keeps ahead of the computation. The
host's peak resident memory is 25.9--26.9~GB.

\paragraph{Optimizer and checkpoints.} Adam ($\beta_1=0.9$, $\beta_2=0.999$, $\epsilon=10^{-8}$, no weight decay, learning
rate $10^{-4}$) updates $B$ from step~1 and $A$ from step~2. After every step the adapters and both moments are written to a
temporary file and renamed; a run of $n$ steps continues from the last checkpoint.

\section{Verification}\label{sec:verify}

\begin{table}[t]\centering\small
\caption{The step is a function of its inputs.}\label{tab:verify}
\begin{tabular}{p{9.2cm}p{5.4cm}}\toprule
test & result \\\midrule
step 1 of the 61-layer model, 1\,024 tokens: 4\,096 vs 40 rows per product, 24 vs 7 threads & 122 $dB$ files byte-identical;
$dB$, $dA$ and $g_0$ digests identical \\
2 steps in one run vs 1 + 1 resumed from the checkpoint (4 layers) & checkpoints byte-identical \\
three generations of product kernels (\texttt{bf16\_gemm2} vs tiled; $[dY|dY][W_{hi};W_{lo}]$ vs E4M3 digits; reference FP8
kernel vs tiled), 4 layers & digests \texttt{63e9904ca4c9d00f} (128 tokens $\times$ 2) and \texttt{1dd8193b18cc200b}
($8\times1024$ + 2 held out) unchanged \\
experts' $dX$, old vs new kernel (layer 3, expert 0, 341 tokens) & 0 outputs differ; 4\,096 checked against the
superaccumulator agree \\
attention backward, GPU vs host (6, 45, 131, 1\,024 tokens) & bit for bit \\
outputs used uncertified & 0 (2--7 a step before the host superaccumulator was added) \\\bottomrule
\end{tabular}
\end{table}

\paragraph{Invariance.} Table~\ref{tab:verify} collects the tests. The third row is the one fixed-order kernels cannot pass:
the three kernel generations differ in tile shapes, number of digits, digit products per term (48 against 18 for the experts'
gradient) and epilogue width (125 against 63 bits), and still give the same bits, because each output is the unique correctly
rounded value. The 63-bit window leaves about 0.02\% of the outputs to the fixup where the 125-bit one left none: the same
outputs from a third of the epilogue's instructions.

\paragraph{Gradients.} Each hand-written backward formula---attention, RoPE, RMSNorm, SwiGLU and the MoE gating---agrees with
central differences of its own FP64 forward on a $2^{-12}$ grid to a relative $5.1\times10^{-7}$, and a deliberately wrong
formula fails. Through the 4-layer model (three dense layers and one MoE layer, 128 tokens), moving the embeddings in 5 random
directions by $\epsilon=0.004$ changes the loss with slope 0.96 against $\langle g_0, \Delta\rangle$; the residual is the
BF16 rounding noise of the forward ($\sim2\times10^{-3}$) and grows with $\epsilon$. On real updates (1\,024 tokens), the loss
changed by $-0.04588$ and $-0.02731$ where the first-order prediction was $-0.04135$ and $-0.03288$.

\section{Performance}\label{sec:perf}

\begin{table}[t]\centering\small
\caption{Training steps of the 61-layer model, $8\times1024$ tokens (+ 2 held-out sequences in the forward), from the training log; ``MoE bwd'' is the MLP and MoE backward, ``attention bwd'' the attention and LoRA backward.}\label{tab:steps}
\begin{tabular}{p{6.2cm}rrrr}\toprule
kernels (step) & step [s] & forward [s] & MoE bwd [s] & attention bwd [s] \\\midrule
digit planes, \texttt{bf16\_gemm2}, GPU attention forward (5)$^{*}$ & 3\,507.9 & 1\,192.6 & 1\,505.8 & 793.1 \\
+ GPU attention recomputed in the backward (7)$^{*}$ & 2\,764.3 & 721.6 & 1\,448.9 & 577.1 \\
+ tiled products, GPU attention backward (9--14) & 1\,967--2\,009 & 634--639 & 884--895 & 462--468 \\
same kernels, at night (16) & 1\,805.5 & & & \\
+ tiled FP8 expert forward (17) & 1\,791.8 & & & \\\bottomrule
\end{tabular}\\[2pt]
{\footnotesize $^{*}$ kernel tests shared the GPU with these steps.}
\end{table}

\begin{table}[t]\centering\small
\caption{Kernels, measured on the laptop while the training shared the GPU.}\label{tab:kernels}
\begin{tabular}{p{9.6cm}r}\toprule
kernel & before $\rightarrow$ after \\\midrule
BF16 product, shared expert's gate, 4\,096 tokens, $p=8$ (TOPS of digit products) & 26.1 $\rightarrow$ 15.7 ms (41 $\rightarrow$ 69) \\
same, $p=24$ & 48.7 $\rightarrow$ 29.2 ms (56 $\rightarrow$ 93) \\
one expert's $dX$ (down + gate/up), 341 tokens, host work included & 81.7 $\rightarrow$ 12.2 ms \\
one expert's forward (gate, $p=8$), 432 tokens (a step's mean) & 1.79 $\rightarrow$ 1.30 ms \\
digit planes of an MoE layer's backward & 27.5 $\rightarrow$ 0.7 s \\
attention backward, 1\,024 tokens, host FP64 $\rightarrow$ GPU FP64 & 2.06 $\rightarrow$ 0.69 s \\\bottomrule
\end{tabular}
\end{table}

A step now takes about 30 minutes (Table~\ref{tab:steps}); the spread between 1\,806 and 2\,009~s with the same kernels is
other work on the same laptop. Table~\ref{tab:kernels} lists the kernel changes behind the fall from 3\,508~s. One of them was
not about arithmetic at all: digit planes written one byte per thread to scattered sectors took 27.5 of the 28~s of an MoE
layer's backward until one warp per block wrote 4-byte words.

\paragraph{What exactness costs.} An exact output needs 12 INT8 products per term for a BF16 result and 30 for an FP32 one. At
69--93~TOPS of digit products that is roughly 15--29 times slower than cuBLAS BF16 on the same GPU (84--91~TFLOPS measured). In the
step this is not the bottleneck: the forward pass's products take 186 of its 524~s. The rest is streaming 391~GB through
host memory twice, the FP64 host work (attention softmax, norms, SwiGLU) and per-expert overheads: 192 experts per layer
are applied one after another, each with its own transfers.

\section{Training on public data}\label{sec:train}

The corpus was harvested from the public websites of Kosovo's institutions and other public sources, politely
(\texttt{robots.txt}, at most one request per second per host). The largest shares come from the Central Bank of the
Republic of Kosovo (40.8 million tokens), the Ministry of Justice (26.0 M), the Ministry of Industry, Entrepreneurship and
Trade (19.3 M) and the Ministry of Education, Science, Technology and Innovation (13.5 M). Lines repeated in at least 20\% of
a source's documents, documents under 500 characters and exact duplicates are removed; Jais-3's tokenizer turns the rest
into 187\,437 sequences of 1\,024 tokens (191.9 million tokens). Each step takes 8 sequences in a seeded order; the last 2
sequences of the corpus are held out and run through every forward.

\begin{table}[t]\centering\small
\caption{Held-out loss (2 sequences, 2\,046 predictions).}\label{tab:loss}
\begin{tabular}{lcccccccc}\toprule
step & 1 & 2 & 3 & 4 & 15 & 16 & 17 & 18 \\\midrule
loss & 1.51019 & 1.50738 & 1.50570 & 1.50056 & 1.46402 & 1.46337 & 1.46709 & 1.46759 \\\bottomrule
\end{tabular}
\end{table}

The held-out loss falls by 0.047 in 16 steps (Table~\ref{tab:loss}) and then holds within the variation of single 8-sequence
batches; after 18 steps the model has seen 147\,456 training tokens, under 0.08\% of the corpus. We report this run to show
that the exact step trains, not as a claim about the adapted model's quality.

\section{Limitations}
\begin{itemize}
\item \emph{Speed.} Thirty minutes per 8\,192 training tokens is slow; the step is a feasibility and reproducibility result
on a laptop, not a throughput result. The per-expert overheads and the FP64 attention backward are the next targets.
\item \emph{What is fixed, and what is unique.} The products are unique; attention, norms and SwiGLU are reproducible
because their FP64 operations run in a fixed order, and their transcendental functions come from the host's C library.
Bit identity across machines therefore assumes the same \texttt{exp}; a correctly rounded library (e.g.\ CORE-MATH) would
remove that assumption.
\item \emph{Modelling choices.} $dY$ is rounded once to BF16 before each transposed product, as BF16 training does; the
top-8 routing is treated as a constant.
\item \emph{Scale of the evaluation.} One model, one adapter configuration, 18 steps; no comparison of downstream quality.
\end{itemize}

\section{Conclusion}
Correct rounding turns a training step into a mathematical function of its inputs. That is a stronger property than
determinism: it survives a change of batch, of thread count, of interruption and of kernel, and it can be checked by anyone
with exact arithmetic, without our code or our hardware. INT8 tensor cores make the products cheap enough that, for a model
streamed from disk, exactness is not what limits the step. We trained a 375-billion-parameter MoE this way on one laptop
GPU, and every one of its products was certified.

\paragraph{Code and data.} The code is source-available under the Ozaki-Blackwell Restricted Source License 1.0 at
\url{https://github.com/ShabanEjupi/ozaki-blackwell}
(files \texttt{llm/jais\_lora.cu} and \texttt{llm/train.sh}); the training logs
print FNV-1a digests of every $dB$ and of the embedding gradient, so a run can be checked bit for bit.

\begin{thebibliography}{99}\small
\bibitem{chen16} T.~Chen, B.~Xu, C.~Zhang, C.~Guestrin. Training deep nets with sublinear memory cost. arXiv:1604.06174, 2016.
\bibitem{exblas} S.~Collange, D.~Defour, S.~Graillat, R.~Iakymchuk. Numerical reproducibility for the parallel reduction on
multi- and many-core architectures. \emph{Parallel Computing} 49:83--97, 2015.
\bibitem{rareblas} C.~Chohra, P.~Langlois, D.~Parello. Reproducible, accurately rounded and efficient BLAS. Euro-Par 2016
Workshops, LNCS 10104, 2017.
\bibitem{reprofuse} L.~Cooper, S.~Jeong, H.~Jeon, J.~Young, H.~Kim. Accelerating the mitigation of LLM inference
nondeterminism across GPU architectures. arXiv:2609.25624, 2026.
\bibitem{dsv3} DeepSeek-AI. DeepSeek-V3 technical report. arXiv:2412.19437, 2024.
\bibitem{dsv4} DeepSeek-AI. DeepSeek-V4: Towards highly efficient million-token context intelligence. arXiv:2606.19348, 2026,
Section~3.3.
\bibitem{reproblas} J.~Demmel, H.~D.~Nguyen. Parallel reproducible summation. \emph{IEEE Trans.\ Computers}
64(7):2060--2070, 2015.
\bibitem{qlora} T.~Dettmers, A.~Pagnoni, A.~Holtzman, L.~Zettlemoyer. QLoRA: Efficient finetuning of quantized LLMs.
arXiv:2305.14314, 2023.
\bibitem{companion} S.~Ejupi. Precision as a software choice: FP64, double-double, triple-double and beyond from the INT8
tensor cores of a laptop GPU. Preprint, 29 October 2026.
\bibitem{he25} H.~He and Thinking Machines Lab. Defeating nondeterminism in LLM inference. Blog post, 10 September 2025,
\url{https://thinkingmachines.ai/blog/defeating-nondeterminism-in-llm-inference/}.
\bibitem{hu21} E.~J.~Hu, Y.~Shen, P.~Wallis, Z.~Allen-Zhu, Y.~Li, S.~Wang, L.~Wang, W.~Chen. LoRA: Low-rank adaptation of
large language models. arXiv:2106.09685, 2021.
\bibitem{kulisch} U.~Kulisch. \emph{Computer Arithmetic and Validity: Theory, Implementation, and Applications}, 2nd ed.
De Gruyter, 2013.
\bibitem{memascend} Y.-C.~Liaw, S.-H.~Chen. MemAscend: System memory optimization for SSD-offloaded LLM fine-tuning.
arXiv:2505.23254, 2025.
\bibitem{ozaki2} K.~Ozaki, Y.~Uchino, T.~Imamura. Ozaki Scheme II: A GEMM-oriented emulation of floating-point matrix
multiplication using an integer modular technique. arXiv:2504.08009, 2025.
\bibitem{zeroinf} S.~Rajbhandari, O.~Ruwase, J.~Rasley, S.~Smith, Y.~He. ZeRO-Infinity: Breaking the GPU memory wall for
extreme scale deep learning. arXiv:2104.07857, 2021.
\bibitem{vllm25} B.~Wasti, W.~Ye, T.~Rao, M.~Goin, P.~Zhang, T.~Liu, N.~Gimelshein, W.~Kwon, K.~You, Z.~Li. No more
train-inference mismatch: bitwise consistent on-policy reinforcement learning with vLLM and TorchTitan. vLLM blog,
10 November 2025, \url{https://vllm.ai/blog/bitwise-consistent-train-inference}.
\bibitem{tbik} Z.~Zhang, X.~Ding, J.~Yuan, R.~Liu, H.~Mao, J.~Xing, Z.~Liu. Deterministic inference across tensor parallel
sizes that eliminates training-inference mismatch. arXiv:2511.17826, 2025.
\bibitem{zhuang21} D.~Zhuang, X.~Zhang, S.~L.~Song, S.~Hooker. Randomness in neural network training: characterizing the
impact of tooling. arXiv:2106.11872, 2021.
\end{thebibliography}
\end{document}
